\section{Conclusion}

In this work, we investigated whether LLMs capture the relational structure of persona-persona dislike. We demonstrated a significant negative correlation between the embedding-based distance of personas and their reciprocal preference agreement. Our results show that LLMs do not merely model isolated traits but possess a sophisticated social map where identity is defined through relational opposition. 

Furthermore, our introduction of the \schismogenesis task revealed that LLMs can effectively invert a target's preferences by constructing a counter-identity, demonstrating a structural understanding of opposition. These findings highlight the inherent "us-vs-them" dynamics in persona-conditioned LLMs, with important implications for the design of cooperative AI and the study of social polarization.

Future work will focus on the asymmetry of dislike and the potential for finding common ground between disparate personas. By understanding how LLMs model social distance, we can better predict and manage their behavior in complex, multi-agent social environments.
